\documentclass[xetex,mathserif,serif]{beamer}

\usepackage{xunicode}
\usepackage{xltxtra}
\usepackage{color}
\usepackage{url}
\usepackage{listings}
\usepackage{fontspec}
\usepackage{geometry}
\usepackage{lastpage}
\usepackage{fancyhdr}
\usepackage{amsmath}
\usepackage{amsthm}
\usepackage{amssymb}
\usepackage{blkarray}
\usepackage{multicol}
\usepackage{relsize}

\definecolor{solarized@base03}{HTML}{002B36}
\definecolor{solarized@base02}{HTML}{073642}
\definecolor{solarized@base01}{HTML}{586e75}
\definecolor{solarized@base00}{HTML}{657b83}
\definecolor{solarized@base0}{HTML}{839496}
\definecolor{solarized@base1}{HTML}{93a1a1}
\definecolor{solarized@base2}{HTML}{EEE8D5}
\definecolor{solarized@base3}{HTML}{FDF6E3}
\definecolor{solarized@yellow}{HTML}{B58900}
\definecolor{solarized@orange}{HTML}{CB4B16}
\definecolor{solarized@red}{HTML}{DC322F}
\definecolor{solarized@magenta}{HTML}{D33682}
\definecolor{solarized@violet}{HTML}{6C71C4}
\definecolor{solarized@blue}{HTML}{268BD2}
\definecolor{solarized@cyan}{HTML}{2AA198}
\definecolor{solarized@green}{HTML}{859900}
\definecolor{yaleblue}{HTML}{0E4C92}

\setbeamertemplate{navigation symbols}{}
% \setbeamerfont{title}{family=\old}
% \setbeamerfont{author}{family=\tfont}%
% \setbeamerfont{frametitle}{family=\oldA}
% \setbeamerfont{date}{family=\dfont}

\setbeamertemplate{itemize items}{--}
\setbeamercolor*{item}{fg=black}

\defaultfontfeatures{Mapping=tex-text}
\hypersetup{pdfstartview={FitH}}

\newcommand{\old}[1]{\fontspec[Alternate=1,Ligatures={Common}]{Hoefler Text}\fontsize{18pt}{30pt}\selectfont #1}%
\newcommand{\oldA}[1]{\fontspec[Alternate=1,Ligatures={Common, Rare}]{Hoefler Text}\fontsize{12pt}{15pt}\selectfont #1}%
\newcommand{\oldB}[1]{\fontspec[Ligatures={Common}]{Didot}\fontsize{12pt}{15pt}\color{solarized@base02}\selectfont #1}%
\newcommand{\tfont}[1]{\fontspec[Alternate=1,Ligatures={Common}]{Hoefler Text}\fontsize{12pt}{20pt}\selectfont #1}%
\newcommand{\dfont}[1]{\fontspec[Ligatures={Common}]{Didot}\fontsize{12pt}{12pt}\selectfont #1}%

\newcommand{\minimize}{\mathop{\mathrm{minimize}}}
\newcommand{\argmin}{\mathop{\mathrm{arg\,min}}}
\newcommand{\argmax}{\mathop{\mathrm{arg\,max}}}
\newcommand{\st}{\mathop{\mathrm{subject\,\,to}}}

\newcommand\independent{\protect\mathpalette{\protect\independenT}{\perp}}
\def\independenT#1#2{\mathrel{\rlap{$#1#2$}\mkern2mu{#1#2}}}

\setlength{\parindent}{0pt}
\setlength{\parskip}{12pt}

\setromanfont [Ligatures={Common}, Numbers={OldStyle}, Variant=01,
 BoldFont={LinLibertine_RB.otf},
 ItalicFont={LinLibertine_RI.otf},
 BoldItalicFont={LinLibertine_RBI.otf}
 ]{LinLibertine_R.otf}



\begin{document}

%%%%%%%%%%%%%%%%%%%%%%%%%%%%%%%%%%%%%%%%%%%%%%%%%%%
\begin{frame}[fragile] \frametitle{}

\vfill

{\fontsize{0.7cm}{0cm}\selectfont Lecture 18 \\\vspace{0.2cm}
Theory of the Lasso}\\\vspace{0.5cm}
30 November 2015

\vspace{2cm}

\begin{minipage}{0.6\textwidth}
Taylor B. Arnold \\
Yale Statistics \\
STAT 312/612
\end{minipage}
\hfill
\begin{minipage}{0.3\textwidth}\raggedleft
\includegraphics[scale=0.3]{../yale-logo.png}
\end{minipage}%

\end{frame}

%%%%%%%%%%%%%%%%%%%%%%%%%%%%%%%%%%%%%%%%%%%%%%%%%%%
\begin{frame}[fragile] \frametitle{}

{\color{yaleblue}\fontsize{16pt}{20pt}\selectfont Class Notes}

\begin{itemize}
\item XXXXX XXXXX XXXXX
\end{itemize}

\end{frame}

%%%%%%%%%%%%%%%%%%%%%%%%%%%%%%%%%%%%%%%%%%%%%%%%%%%
\begin{frame}[fragile] \frametitle{}

Recall that sasso regression replaces the $\ell_2$ penalty with an $\ell_1$ penalty,
and looks deceptively similar to the ridge regression:
\begin{align*}
\widehat{\beta}_\lambda &= \argmin_b \left\{ || y - Xb ||_2^2 + \lambda ||b||_1 \right\}
\end{align*}
Where the $\ell_1$-norm is defined as the sum of the absolute values of the
vector's components:
\begin{align*}
|| \beta ||_1 &= \sum_i | \beta_i |
\end{align*}

\end{frame}

%%%%%%%%%%%%%%%%%%%%%%%%%%%%%%%%%%%%%%%%%%%%%%%%%%%
\begin{frame}[fragile] \frametitle{}

There are three types of errors that we commonly are concerned
with in lasso regression. \pause The prediction loss:
\begin{align*}
|| X (\beta - \widehat{\beta}) ||_2^2
\end{align*}
\pause Parameter estimation:
\begin{align*}
|| \beta - \widehat{\beta} ||_2^2
\end{align*}
\pause And model selection:
\begin{align*}
\text{supp} (\beta) &= \text{supp} (\widehat{\beta})
\end{align*}

\pause Today we are going to focus on the first two types
of estimation errors.

\end{frame}

%%%%%%%%%%%%%%%%%%%%%%%%%%%%%%%%%%%%%%%%%%%%%%%%%%%
\begin{frame}[fragile] \frametitle{}

The type of results we want to establish may look different than
you are used to seeing in more introductory courses.
\pause What we will want to be able to construct is a set
$\mathcal{A}$ such that:
\begin{align*}
\mathbb{P} \mathcal{A} &= 1 - \epsilon
\end{align*}
For some small $\epsilon > 0$, and
\begin{align*}
|| X (\beta - \widehat{\beta}) ||_2^2 &\leq \delta
\end{align*}
Conditioned on being in event $\mathcal{A}$.

\end{frame}

%%%%%%%%%%%%%%%%%%%%%%%%%%%%%%%%%%%%%%%%%%%%%%%%%%%
\begin{frame}[fragile] \frametitle{}

We want to establish some result about the lasso solution; where
to begin?

Let's start with the one relationship we know to be true between
$\widehat{\beta}$ and $\beta$:
\begin{align*}
|| y - X\widehat{\beta} ||_2^2 + \lambda ||\widehat{\beta}||_1 \leq || y - X\beta ||_2^2 + \lambda ||\beta||_1
\end{align*}

\end{frame}

%%%%%%%%%%%%%%%%%%%%%%%%%%%%%%%%%%%%%%%%%%%%%%%%%%%
\begin{frame}[fragile] \frametitle{}

I'll start by expanding the loss terms and canceling on both sides:
\begin{align*}
|| y - X\widehat{\beta} ||_2^2 + \lambda ||\widehat{\beta}||_1 &\leq || y - X\beta ||_2^2 + \lambda ||\beta||_1 \\
y^t y + \widehat{\beta}^t X^t X \widehat{\beta} - 2 y^t X \widehat{\beta} + \lambda ||\widehat{\beta}||_1
&\leq y^t y + b^t X^t X b - 2 y^t X b + \lambda ||b||_1\\
\widehat{\beta}^t X^t X \widehat{\beta} - 2 y^t X \widehat{\beta} + \lambda ||\widehat{\beta}||_1
&\leq b^t X^t X b - 2 y^t X b + \lambda ||b||_1\\
\end{align*}

\end{frame}

%%%%%%%%%%%%%%%%%%%%%%%%%%%%%%%%%%%%%%%%%%%%%%%%%%%
\begin{frame}[fragile] \frametitle{}

The next step is grouping the terms on each side of the equation
\begin{align*}
\widehat{\beta}^t X^t X \widehat{\beta} - 2 y^t X \widehat{\beta} + \lambda ||\widehat{\beta}||_1
&\leq b^t X^t X b - 2 y^t X b + \lambda ||b||_1\\
|| X (\widehat{\beta} - b) ||_2^2 + \lambda ||\widehat{\beta}||_1
&\leq 2 y^t X (\beta - \widehat{\beta}) + \lambda ||b||_1
\end{align*}

\end{frame}


\end{document}











